% Section 7: Evaluation
% Target length: 800-1500 words
% Status: PILOT RUN 1 COMPLETE (Chapter 12, N=1 chapter, LLM-judge only --
% see "Pilot Results" subsection below). This is a first, small, promising
% result, NOT the powered 6-10-chapter study with independent human raters
% that the Design subsection still calls for. Do not let the pilot's positive
% result relax the standard for the full study -- see caveats inline.

\section{Evaluation}
\label{sec:evaluation}

\subsection{Design: Blind Comparison Study}
% TODO: expand into prose.
\begin{itemize}
    \item \textbf{Critical control:} hold total chapter word count constant across
          conditions. Each condition redistributes the same budget differently;
          this isolates the allocation effect from a simple ``more words = better''
          confound.
    \item \textbf{Conditions} (within-chapter, same concepts, same model = Claude
          Sonnet 5, same prompt template except the per-concept word target):
          \begin{itemize}
              \item[A --] Uniform: $\mathrm{words}(c) = W_i / |Ch_i|$ for all $c$.
              \item[B --] In-degree weighted (specified below).
              \item[C --] CIS weighted (proposed): see \S\ref{sec:problem-statement}.
          \end{itemize}
    \item \textbf{Pilot corpus:} Algebra I learning graph and chapters
          (\texttt{/Users/danmccreary/Documents/ws/algebra-1}). Chapter 12
          (``Exponential Functions'', 10 concepts) selected for the first pilot
          run -- already validated as the specialized-chapter contrast case
          (\S\ref{sec:system}), and its small size (3{,}722 published words)
          keeps the first full 3-condition run tractable. Chapter 1 remains
          the planned second pilot chapter (largest CIS/in-degree divergence,
          \S\ref{sec:system}).
    \item \textbf{Sample:} 6-10 chapters total, spanning varied CIS distributions
          (some with one dominant hub concept, some flatter) and, if feasible, more
          than one book/domain so the finding is not a single-domain artifact.
\end{itemize}

\subsection{Condition B Specification: In-Degree Weighted}
\label{sec:condition-b}

There is no single documented ``existing in-degree formula'' to reproduce --
the original observation (\S\ref{sec:introduction}) came from an informal
prompting session (wake-word-detection project, Google Antigravity) whose
output was not preserved. Condition B is therefore specified fresh, using
the \emph{same functional form} as Condition C so that the only variable
between B and C is which importance metric drives the weights -- this is
what makes the comparison a fair test of CIS specifically, not just ``any
graph-derived weighting beats uniform.''

\[
E_B(c) \;=\; \frac{\log(\mathrm{indeg}(c) + 1)}{\log(\mathrm{indeg}_{\max} + 1)}, \qquad
\mathrm{words}_B(c) \;=\; w_{\min} + (W_i - n \cdot w_{\min}) \cdot \frac{\log(\mathrm{indeg}(c)+1)}{\sum_{c' \in Ch_i} \log(\mathrm{indeg}(c')+1)}
\]

with $\mathrm{indeg}_{\max}$ the maximum in-degree over the whole book (16,
``Function'', in the Algebra~I graph) and $w_{\min}=150$, matching Condition
C's floor. Unlike CIS, raw in-degree can be exactly 0 (a concept nobody
directly depends on), which log-normalizes to $E_B=0$ cleanly -- no
floor-tie pathology analogous to Definition 4's CIS-population-rank bug,
since $\log(0+1)=0$ is already the natural bottom of the scale rather than a
population-dependent cutoff.

\begin{table}[htbp]
\centering
\caption{All three conditions' word budgets for Algebra I, Chapter 12, fixed 3{,}722-word total. Condition B is visibly more extreme than C: every concept with in-degree 0 is pinned to the 150-word floor, while CIS's transitive smoothing still differentiates among them (Table~\ref{tab:ch12-budget}).}
\label{tab:ch12-three-conditions}
\begin{tabular}{lrr rr r}
\toprule
Concept & in-deg & CIS & A (unif.) & B (in-deg) & C (CIS) \\
\midrule
Exponential Function             & 6 & 10 & 372 & 1307 & 697 \\
Exponential Growth                & 2 & 3  & 372 & 803  & 466 \\
Exponential Decay                 & 1 & 2  & 372 & 562  & 401 \\
Growth Factor                     & 0 & 1  & 372 & 150  & 308 \\
Decay Factor                      & 0 & 1  & 372 & 150  & 308 \\
Initial Value                     & 0 & 1  & 372 & 150  & 308 \\
Exponential Models                & 0 & 1  & 372 & 150  & 308 \\
Comparing Linear and Exponential  & 0 & 1  & 372 & 150  & 308 \\
Graphing Exponentials             & 0 & 1  & 372 & 150  & 308 \\
Applications of Systems           & 0 & 1  & 372 & 150  & 308 \\
\bottomrule
\end{tabular}
\end{table}

Condition B allocates \emph{more than three times} the budget to
``Exponential Function'' (1307 words) as Condition C (697), and slams all
seven zero-in-degree concepts to the bare 150-word floor rather than C's
still-differentiated 308. This is the concrete, testable form of the paper's
central claim: raw in-degree treats every concept nothing directly depends
on as equally, minimally important, while CIS's transitive reach still
distinguishes among them by how much the book eventually builds on the
surrounding neighborhood -- even when direct dependents are zero.

\subsection{Rating Protocol}
% TODO: expand.
\begin{itemize}
    \item \textbf{LLM-judge rubric:} separate model call, blinded to condition,
          scoring 1-5 on: (a) depth/elaboration on foundational concepts,
          (b) use of examples/metaphors/contrasting cases, (c) coherence,
          (d) absence of padding on minor concepts. Run 2-3 times per chapter to
          check judge consistency (report agreement, e.g. intraclass correlation).
    \item \textbf{Human spot-check:} a subset of chapters (e.g. 3 of 6-10),
          blind-rated by the author and, if available, 1-2 additional raters, same
          rubric, condition labels stripped, presentation order randomized.
          Purpose: validate the LLM-judge against human judgment and report their
          agreement -- this is itself a result worth including.
\end{itemize}

\subsection{Analysis Plan}
% TODO: expand.
\begin{itemize}
    \item Small $N$ expected -- do not assume normality. Report pairwise win/tie/loss
          tallies (C vs.\ A, C vs.\ B, B vs.\ A) with a sign test or Wilcoxon
          signed-rank test.
    \item Secondary check (if feasible): correlate per-concept word allocation
          under each condition against an independent difficulty/importance
          signal (e.g. instructor judgment or exam-weighting data, if available)
          as a construct-validity check on CIS itself, not just on downstream
          content quality.
\end{itemize}

\subsection{Pilot Results: Chapter 12}
\label{sec:pilot-ch12}

% TODO: expand into prose. Full generation transcript, raw judge output, and
% blinding methodology recorded in pilot-study/chapter-12/RESULTS.md for
% reproducibility.

\textbf{Generation.} Three complete chapter drafts were generated
independently -- separate sessions, no shared context, none with access to
the others' output or to the real published Chapter 12 -- from the same
fixed scaffold (title, summary, prerequisites, the 10 concepts in order),
differing only in the per-concept word-budget instruction
(Table~\ref{tab:ch12-three-conditions}). All three landed within the
$\pm15\%$ per-concept tolerance and within 1\% of the 3{,}722-word chapter
total (achieved: A 3{,}626; B 3{,}638; C 3{,}736).

\textbf{Blind judging.} Word Count Report footers (which would leak the
condition via the target numbers) were stripped, files were relabeled
Draft~1/2/3 in an order not disclosed to the judges, and three independent
fresh model sessions each scored all three drafts on the four-dimension
rubric (\S\ref{sec:evaluation}, Rating Protocol) without seeing each other's
scores or this paper's hypothesis.

\begin{table}[htbp]
\centering
\caption{Pilot judge scores (1-5, mean of 3 independent blind judges), Chapter 12.}
\label{tab:pilot-ch12-scores}
\begin{tabular}{lrrrrr}
\toprule
Condition & a: depth & b: examples & c: coherence & d: no padding & overall \\
\midrule
A (uniform)         & 3.00 & 3.33 & 3.67 & 3.00 & 3.00 \\
B (in-degree)        & 4.33 & 4.00 & 4.00 & 4.67 & 4.33 \\
C (CIS, proposed)    & 4.67 & 5.00 & 5.00 & 4.00 & 5.00 \\
\bottomrule
\end{tabular}
\end{table}

\textbf{Result.} All three judges independently produced the identical
ranking: C $>$ B $>$ A (3/3 agreement). Notably B also beat A in this pilot --
even in-degree's cruder, more extreme skew (Table~\ref{tab:ch12-three-conditions})
outperformed uniform allocation here -- with C ahead of both.

\textbf{Qualitative illustration.} On ``Growth Factor'' (a low-budget concept
under both B and C), Condition B (154 words) produces a single definitional
paragraph with no worked example; Condition C (290 words), despite a similar
target, still fits a full labeled worked example (a concrete population-growth
scenario) plus a second technique (recovering the growth factor from a data
table). This is the mechanism the aggregate scores above are picking up on.

\textbf{Caveats -- read before citing this result.} This is $N=1$ chapter,
not the powered 6-10-chapter sample the Design subsection calls for. The
three ``blind judges'' are fresh sessions of the same model family used to
generate the drafts (Claude Sonnet 5), not independent human raters --
same-model self-consistency cannot be ruled out as a contributor to the
result, and the planned human spot-check (\S\ref{sec:evaluation}, Rating
Protocol) has not yet been run. Chapter 12 is a specialized, leaf-heavy
chapter (\S\ref{sec:system}, Table~\ref{tab:ch12-budget}); the equivalent run
on foundational Chapter 1 (opposite CIS profile) is still outstanding. Treat
this as a promising first pilot signal that justifies running the full
study, not as the study's result.

\subsection{Resolved Items}
\begin{itemize}
    \item \textbf{Global vs. local CIS normalization -- RESOLVED: global.} See
          \S\ref{sec:formal-definitions} and Table~\ref{tab:ch1-budget}
          (\S\ref{sec:system}). Condition C content generation should
          normalize against the full 200-concept Algebra I graph, not just
          the concepts in one chapter.
    \item \textbf{Tier-threshold calibration on a non-foundational chapter --
          RESOLVED, with a correction along the way.} Computing the same
          table for Chapter 12 exposed that the original tiering rule
          (population percentile rank) was structurally broken: CIS has a
          floor of 1, 103/200 Algebra I concepts (51.5\%) sit at that floor,
          and the population 35th-percentile cutoff landed \emph{at} the
          floor -- making Tier C unreachable regardless of true importance.
          Fixed by switching to the log-scaled Elaboration Score $E(c)$
          (Definition 4, \S\ref{sec:formal-definitions}), which separates
          floor-tied concepts from true hubs. Re-validated on both chapters:
          Chapter 1 now gives a real 8/6/3 A/B/C split (previously 11/6/0
          under percentile rank); Chapter 12 gives 0/2/8, the expected
          mirror-image profile for a late, specialized chapter
          (Tables~\ref{tab:ch1-budget}-\ref{tab:ch12-budget}). Condition C
          generation should use $E(c)$, not population percentile.
\end{itemize}

\subsection{Open Items Before Completing the Study}
% TODO: resolve each of these, then delete this subsection.
\begin{itemize}
    \item \textbf{Independent human rater(s) for Chapter 12} -- the pilot
          result (\S\ref{sec:pilot-ch12}) currently rests entirely on
          same-model-family LLM judges. This is the single most important
          remaining gap before the result can be reported as more than
          preliminary.
    \item \textbf{Run the same 3-condition generation + blind judging on
          Chapter 1} (foundational profile, opposite of Chapter 12's
          specialized profile) to check the pattern isn't specific to
          leaf-heavy chapters.
    \item \textbf{Extend to the full 6-10 chapter sample} if the Chapter 1
          result is consistent with Chapter 12's.
    \item \textbf{$E(c)$ cut points (0.5/0.2) are only lightly validated.}
          Confirmed sensible on two chapters with opposite expected profiles,
          but that is a small validation set.
    \item \textbf{Elaboration-budget dimensions beyond words} (diagrams,
          MicroSims, worked examples) -- the pilot varied word count only;
          decide whether to extend to the full multi-dimensional elaboration
          budget in later runs.
\end{itemize}

% TODO: once the study runs, replace this section's outline with: study
% execution details, results table/figure, statistical results, and
% interpretation.
